% [DISCURSO D04] Defina desde aqui o referente de «antropocênicos»: a humanidade
%   europeia-modernizante-ocidental que se toma pela espécie. Em notas/antropocenicos.tex cabe
%   explicitar que se trata de um modo histórico de compor coletivos, cuja adesão pode variar; a
%   palavra não designa uma origem geográfica ou todos os europeus.
Podemos imaginar os antropocênicos como aqueles acometidos pela cosmofobia, adoecidos em sua comunicação e ancestralidade. A patologia se articula à humanidade europeia-modernizante-ocidental, embora possa acometer qualquer indivíduo, e será melhor descrita adiante. O \enquote{Antropoceno} é uma nomeação \textit{pelos} antropocênicos, e o uso das aspas indica, entre outras coisas, que falamos enquanto não-antropocênicos.